% GENERATED FILE. Edit canonical lessons, then run npm run generate:books.
% canonical-source-hash: fnv1a64:f1b5ab0fb0e3ce41
% canonical-lessons: UR-C321-jhulna, UR-C321-julus, UR-C321-hujum, UR-C321-ikattha, UR-C321-tamasha

\chapter{At the Fair}
\label{ch:ur-at-the-fair}
\hlchaptermodality{321}

\begin{chapteropening}
\textbf{By the end of this chapter:} \emph{I can describe the crowds and the sights at a fair.}

Complete the last lesson of chapter 321: I can describe the crowds and the sights at a fair.
\end{chapteropening}

\section[\texorpdfstring{\ur{جھولنا} (jhūlnā)}{جھولنا (jhūlnā)}]{\ur{جھولنا} (jhūlnā) — to swing}

\label{lesson:UR-C321-jhulna}



\begin{quote}
Before the new one: say the Urdu for to drench, then the Urdu for to sway.
\end{quote}

\subsection*{You'll want to know: \ur{جھولنا}}
\textbf{\ur{جھولنا}} — \emph{jhūlnā} — \textquotedblleft{}to swing\textquotedblright{}.

\begin{quote}
\ur{بچی} \ur{رسی} \ur{پر} \ur{جھول} \ur{رہی} \ur{تھی۔}
\end{quote}

\begin{quote}
\emph{bachchī rassī par jhūl rahī thī.}
\end{quote}

\begin{quote}
The little girl was swinging on a rope.
\end{quote}

Past (she …): \ur{وہ} \ur{جھولی}, \emph{voh jhūlī}.

\begin{cousinweb}
\textbf{\ur{جھولنا}} (\emph{jhūlnā}) is inherited: CDIAL lists the Hindi cousin \emph{jhūlnā}, \textquotedblleft{}to swing\textquotedblright{}, under a form it reconstructs (CDIAL 5406).
\end{cousinweb}

\begin{grammarlens}[title={What you've built}]
One more word for telling what happened.
\end{grammarlens}

\subsection*{Your turn}
\begin{itemize}
\raggedright
  \item \emph{Say it:} \emph{jhūlnā}
  \item \emph{Say it:} \emph{jhūlnā}, once more
  \item \emph{Say it:} \emph{bachchī rassī par jhūl rahī thī.} — The little girl was swinging on a rope.
  \item \emph{Recall:} say the Urdu for to drench, then the Urdu for to sway, then say \emph{jhūlnā} again
\end{itemize}

\begin{tcolorbox}[breakable,colback=teal!4,colframe=teal!35!black,title={Before you move on}]
What does \ur{جھولنا} mean? (\textquotedblleft{}To swing\textquotedblright{}.) Say it once more.
\end{tcolorbox}
\section[\texorpdfstring{\ur{جلوس} (julūs)}{جلوس (julūs)}]{\ur{جلوس} (julūs) — a procession}

\label{lesson:UR-C321-julus}



\begin{quote}
Before the new one: say the Urdu for to sway, then the Urdu for to swing.
\end{quote}

\subsection*{You'll want to know: \ur{جلوس}}
\textbf{\ur{جلوس}} — \emph{julūs} — \textquotedblleft{}a procession\textquotedblright{}.

\begin{quote}
\ur{سڑک} \ur{پر} \ur{ایک} \ur{لمبا} \ur{جلوس} \ur{تھا۔}
\end{quote}

\begin{quote}
\emph{saṛak par ek lambā julūs thā.}
\end{quote}

\begin{quote}
There was a long procession in the road.
\end{quote}

Masculine.

\begin{cousinweb}
\textbf{\ur{جلوس}} (\emph{julūs}) is usually traced, through Persian, to Arabic \emph{julūs}, \textquotedblleft{}a sitting\textquotedblright{}, once the sitting of a ruler in state; how it came to mean a procession is thought to run through the royal parades.
\end{cousinweb}

\begin{grammarlens}[title={What you've built}]
One more word for telling what happened.
\end{grammarlens}

\subsection*{Your turn}
\begin{itemize}
\raggedright
  \item \emph{Say it:} \emph{julūs}
  \item \emph{Say it:} \emph{julūs}, once more
  \item \emph{Say it:} \emph{saṛak par ek lambā julūs thā.} — There was a long procession in the road.
  \item \emph{Recall:} say the Urdu for to sway, then the Urdu for to swing, then say \emph{julūs} again
\end{itemize}

\begin{tcolorbox}[breakable,colback=teal!4,colframe=teal!35!black,title={Before you move on}]
What does \ur{جلوس} mean? (\textquotedblleft{}A procession\textquotedblright{}.) Say it once more.
\end{tcolorbox}
\section[\texorpdfstring{\ur{ہجوم} (hujūm)}{ہجوم (hujūm)}]{\ur{ہجوم} (hujūm) — a crowd}

\label{lesson:UR-C321-hujum}



\begin{quote}
Before the new one: say the Urdu for to swing, then the Urdu for a procession.
\end{quote}

\subsection*{You'll want to know: \ur{ہجوم}}
\textbf{\ur{ہجوم}} — \emph{hujūm} — \textquotedblleft{}a crowd\textquotedblright{}.

\begin{quote}
\ur{میلے} \ur{میں} \ur{بڑا} \ur{ہجوم} \ur{تھا۔}
\end{quote}

\begin{quote}
\emph{mele meṅ baṛā hujūm thā.}
\end{quote}

\begin{quote}
There was a big crowd at the fair.
\end{quote}

Masculine.

\begin{cousinweb}
\textbf{\ur{ہجوم}} (\emph{hujūm}) is usually traced, through Persian, to Arabic \emph{hujūm}, \textquotedblleft{}a rushing in, an onset\textquotedblright{}; in Urdu it is the crowd itself.
\end{cousinweb}

\begin{grammarlens}[title={What you've built}]
One more word for telling what happened.
\end{grammarlens}

\subsection*{Your turn}
\begin{itemize}
\raggedright
  \item \emph{Say it:} \emph{hujūm}
  \item \emph{Say it:} \emph{hujūm}, once more
  \item \emph{Say it:} \emph{mele meṅ baṛā hujūm thā.} — There was a big crowd at the fair.
  \item \emph{Recall:} say the Urdu for to swing, then the Urdu for a procession, then say \emph{hujūm} again
\end{itemize}

\begin{tcolorbox}[breakable,colback=teal!4,colframe=teal!35!black,title={Before you move on}]
What does \ur{ہجوم} mean? (\textquotedblleft{}A crowd\textquotedblright{}.) Say it once more.
\end{tcolorbox}
\section[\texorpdfstring{\ur{اکٹھا} (ikaṭṭhā)}{اکٹھا (ikaṭṭhā)}]{\ur{اکٹھا} (ikaṭṭhā) — gathered together}

\label{lesson:UR-C321-ikattha}



\begin{quote}
Before the new one: say the Urdu for a procession, then the Urdu for a crowd.
\end{quote}

\subsection*{You'll want to know: \ur{اکٹھا}}
\textbf{\ur{اکٹھا}} — \emph{ikaṭṭhā} — \textquotedblleft{}gathered together\textquotedblright{}.

\begin{quote}
\ur{سب} \ur{لوگ} \ur{میدان} \ur{میں} \ur{اکٹھے} \ur{ہوئے۔}
\end{quote}

\begin{quote}
\emph{sab log maidān meṅ ikaṭṭhe hue.}
\end{quote}

\begin{quote}
Everyone gathered on the field.
\end{quote}

\begin{cousinweb}
\textbf{\ur{اکٹھا}} (\emph{ikaṭṭhā}) is inherited: CDIAL lists the Hindi cousin \emph{ekṭhā}, \textquotedblleft{}standing together\textquotedblright{}, under Sanskrit \emph{ēkasthá} (CDIAL 2480). It agrees like an adjective: \textbf{\ur{اکٹھے} \ur{ہوئے}}.
\end{cousinweb}

\begin{grammarlens}[title={What you've built}]
One more word for telling what happened.
\end{grammarlens}

\subsection*{Your turn}
\begin{itemize}
\raggedright
  \item \emph{Say it:} \emph{ikaṭṭhā}
  \item \emph{Say it:} \emph{ikaṭṭhā}, once more
  \item \emph{Say it:} \emph{sab log maidān meṅ ikaṭṭhe hue.} — Everyone gathered on the field.
  \item \emph{Recall:} say the Urdu for a procession, then the Urdu for a crowd, then say \emph{ikaṭṭhā} again
\end{itemize}

\begin{tcolorbox}[breakable,colback=teal!4,colframe=teal!35!black,title={Before you move on}]
What does \ur{اکٹھا} mean? (\textquotedblleft{}Gathered together\textquotedblright{}.) Say it once more.
\end{tcolorbox}
\section[\texorpdfstring{\ur{تماشا} (tamāshā)}{تماشا (tamāshā)}]{\ur{تماشا} (tamāshā) — a show; a spectacle}

\label{lesson:UR-C321-tamasha}



\begin{quote}
Before the new one: say the Urdu for a crowd, then the Urdu for gathered together.
\end{quote}

\subsection*{You'll want to know: \ur{تماشا}}
\textbf{\ur{تماشا}} — \emph{tamāshā} — \textquotedblleft{}a show; a spectacle\textquotedblright{}.

\begin{quote}
\ur{میلے} \ur{میں} \ur{بندر} \ur{کا} \ur{تماشا} \ur{تھا۔}
\end{quote}

\begin{quote}
\emph{mele meṅ bandar kā tamāshā thā.}
\end{quote}

\begin{quote}
There was a monkey show at the fair.
\end{quote}

Masculine.

\begin{cousinweb}
\textbf{\ur{تماشا}} (\emph{tamāshā}) is usually traced, through Persian, to Arabic \emph{tamāshā}, \textquotedblleft{}a walking about together\textquotedblright{}, and so a sight people gather to watch.
\end{cousinweb}

\begin{grammarlens}[title={What you've built}]
One more word for telling what happened.
\end{grammarlens}

\subsection*{Your turn}
\begin{itemize}
\raggedright
  \item \emph{Say it:} \emph{tamāshā}
  \item \emph{Say it:} \emph{tamāshā}, once more
  \item \emph{Say it:} \emph{mele meṅ bandar kā tamāshā thā.} — There was a monkey show at the fair.
  \item \emph{Recall:} say the Urdu for a crowd, then the Urdu for gathered together, then say \emph{tamāshā} again
\end{itemize}

\begin{tcolorbox}[breakable,colback=teal!4,colframe=teal!35!black,title={Before you move on}]
What does \ur{تماشا} mean? (\textquotedblleft{}A show; a spectacle\textquotedblright{}.) Say it once more.
\end{tcolorbox}
